\section{Introduction}
\label{sec:intro}

Parton distribution functions (PDFs) play a fundamental role in high-energy particle physics, encoding the momentum distributions of quarks and gluons within protons and neutrons.
These functions are essential inputs for making precise predictions at hadron colliders such as the Large Hadron Collider (LHC), directly impacting the interpretation of a wide range of Standard Model (SM) and beyond-the-SM (BSM) processes.
Among the PDFs, the gluon distribution is particularly crucial due to the dominant role gluons play in the proton's momentum at high energies and in processes like Higgs-boson production and jet formation.
While quark PDFs are relatively well constrained by deep inelastic scattering and Drell-Yan data, the gluon PDF remains less precisely determined, especially in the intermediate and large momentum fraction ($x$) regions.
This uncertainty limits the precision of theoretical predictions for key observables at the LHC.
Improving the accuracy of the gluon PDF, therefore, has the potential to significantly enhance the sensitivity of searches for BSM physics and refine our understanding of the dynamics of quantum-chromodynamics (QCD) in hadronic collisions~\cite{Amoroso:2022eow}.

Lattice QCD (LQCD) provides a first-principles approach to studying the nonperturbative regime of the strong interaction by discretizing spacetime on a finite lattice.
Although LQCD was traditionally limited to the calculation of a few hadron structure observables such as moments of PDFs, recent breakthroughs have enabled LQCD calculations to access the Bjorken-$x$--dependence of PDFs.
Methods such as Large-Momentum Effective Theory (LaMET)~\cite{Ji:2013dva} and pseudo-PDFs~\cite{Radyushkin:2017cyf} are the most popular approaches to extracting lightcone correlations from lattice observables.
Over the past decade or so, substantial progress has been made not only in PDFs but also in other kinds of three-dimensional structures, such as generalized parton distributions (GPDs);
we refer interested readers to the reviews in Refs.~\cite{Ji:2020ect,Constantinou:2020hdm,Lin:2023kxn}.
However, the majority of lattice $x$-dependent results focus on the isovector quark distributions, and there are few works on the gluon PDFs.

The extraction of gluon PDFs in lattice QCD remains more challenging, due to their inherently noisier signals and the need to handle more complicated operator mixing.
Nevertheless, encouraging progress has been reported in recent years, mostly using the pseudo-PDF framework, obtaining the nucleon~\cite{Fan:2020cpa, Fan:2022kcb, HadStruc:2021wmh, Delmar:2023agv, HadStruc:2022yaw}, pion~\cite{Good:2023ecp}, and kaon~\cite{Salas-Chavira:2021wui} gluon PDFs.
However, the pseudo-PDF method often requires the use of phenomenology-inspired $x$-dependent PDF models to fit the ratios of lattice matrix elements; though, there are recent efforts to reduce model dependence in this methodology \cite{Karpie:2019eiq, Karpie:2021pap}.
LaMET, on the other hand, uses a Fourier transform of the lattice matrix elements to first obtain quasi-PDFs, which can be matched to the lightcone PDFs via a matching procedure accurate to powers of the inverse parton momentum.
Model dependence comes into play in LaMET in the extrapolation of large-distance matrix elements to obtain a reliable Fourier transform.
It would be of great community interest to obtain the gluon PDF through LaMET as well, so that we can achieve better control of the lattice systematics by comparing results from the two methodologies.

Recently, MSULat group aimed to advance the LQCD calculations of gluon PDFs using LaMET~\cite{Good:2024iur}.
We used high-statistics measurements at two pion masses (approximately 310 and 690~MeV) on a lattice ensemble with a spacing of about 0.12~fm.
The work explored the application of Coulomb gauge fixing to enhance signal quality at large Wilson-line separations, demonstrating its potential to improve the reliability of the calculations of gluon matrix elements.
By analyzing multiple gluon operators and comparing the resulting matrix elements for both nucleon and pion states, the study identified the most effective operator for extracting gluon PDFs.
However, without the perturbative QCD (pQCD) calculations needed for hybrid-ratio renormalization of this operator, we could only perform an analysis with an estimation of the renormalization.

In this work, we took the lattice gluon matrix elements from our previous work~\cite{Good:2024iur} and calculated the hybrid-renormalization-scheme matching coefficients to obtain the first lightcone gluon PDF of the nucleon using the LaMET approach.
We report the next-to-leading-order (NLO) Wilson coefficients for the hybrid renormalization scheme for the gluon operator in Sec.~\ref{sec:analytic}.
We determine the counterterms in the hybrid renormalization and apply them to renormalize the LaMET matrix elements in Sec.~\ref{sec:numerical}.
We compare the renormalized matrix elements with those estimated from phenomenology, and present the first nucleon gluon PDF from LaMET, also compared with phenomenological fits.
We conclude and consider the future prospects for gluon PDFs in Sec.~\ref{sec:conclusion}.
